\subsection{Injective modules}

DEFINITION 2. --- The A-module E is said to be injective if, for every exact sequence of A-modules and homomorphisms

\[
\mathbf {M} ^ {\prime} \xrightarrow {u} \mathbf {M} \xrightarrow {v} \mathbf {M} ^ {\prime \prime},\tag{19}
\]

the sequence of Z-linear maps

\[
\operatorname{Hom} _ {\mathbf {A}} \left(\mathbf {M} ^ {\prime \prime}, \mathbf {E}\right) \xrightarrow {\operatorname{Hom} _ {\mathbf {A}} (v , 1)} \operatorname{Hom} _ {\mathbf {A}} (\mathbf {M}, \mathbf {E}) \xrightarrow {\operatorname{Hom} _ {\mathbf {A}} (u , 1)} \operatorname{Hom} _ {\mathbf {A}} \left(\mathbf {M} ^ {\prime}, \mathbf {E}\right)\tag{20}
\]

is exact.

Lemma 3. --- For the A-module E to be injective, it is necessary and sufficient that, for every injective A-linear map u : M' → M, the map

\[
\operatorname{Hom} _ {\mathbf {A}} (u, 1): \operatorname{Hom} _ {\mathbf {A}} (\mathrm{M}, \mathrm{E}) \rightarrow \operatorname{Hom} _ {\mathbf {A}} (\mathrm{M} ^ {\prime}, \mathrm{E})
\]

be surjective.

If E is injective and if \(u: \mathbf{M}' \to \mathbf{M}\) is injective, then the sequence \(0 \to \mathbf{M}' \xrightarrow{u} \mathbf{M}\) is exact, hence so is the sequence Hom (M, E) \(\xrightarrow{\text{Hom}(u,1)}\) Hom ( \(\mathbf{M}'\) , E) \(\to 0\) , and Hom ( \(u,1\) ) is surjective. Conversely, consider the exact sequence (19); set \(\mathbf{M}_{1}'' = v(\mathbf{M})\) and let \(i: \mathbf{M}_{1}'' \to \mathbf{M}''\) be the canonical injection and \(p: \mathbf{M} \to \mathbf{M}_{1}''\) the map \(m \mapsto v(m)\) . The sequence \(\mathbf{M}' \xrightarrow{u} \mathbf{M} \xrightarrow{p} \mathbf{M}_{1}'' \to 0\) is exact; by II, p.~36, th. 1, the sequence

\[
\operatorname{Hom} _ {\mathbf {A}} \left(\mathrm{M} _ {1} ^ {\prime \prime}, \mathrm{E}\right) \xrightarrow {\operatorname{Hom} (p , 1)} \operatorname{Hom} _ {\mathbf {A}} (\mathrm{M}, \mathrm{E}) \xrightarrow {\operatorname{Hom} (u , 1)} \operatorname{Hom} _ {\mathbf {A}} \left(\mathrm{M} ^ {\prime}, \mathrm{E}\right)
\]

is exact. Moreover, we have Hom \((v, 1) = \text{Hom } (p, 1) \circ \text{Hom } (i, 1)\) . If E satisfies the condition of the lemma, Hom \((i, 1)\) is surjective, so the image of Hom \((v, 1)\) is also that of Hom \((p, 1)\) , and the sequence (20) is exact.

Remark. --- Let E be an injective A-module, \(u: \mathbf{M}' \to \mathbf{M}\) and \(f: \mathbf{M}' \to \mathbf{E}\) homomorphisms of A-modules. If \(\operatorname{Ker} u \subset \operatorname{Ker} f\) , there exists a homomorphism \(g: \mathbf{M} \to \mathbf{E}\) such that \(g \circ u = f\) . This follows in fact from what precedes, applied to the injective homomorphism \(\mathbf{M}' / \operatorname{Ker} u \to \mathbf{M}\) deduced from \(u\) .

PROPOSITION 9. --- Let \((\mathbf{E}_i)_{i\in I}\) be a family of A-modules, \(\mathbf{E} = \prod \mathbf{E}_i\) their product. For the A-module E to be injective, it is necessary and sufficient that each of the \(\mathbf{E}_i\) be so.

Let \(u: \mathbf{M}' \to \mathbf{M}\) be an injective homomorphism of A-modules. For the product homomorphism \(\prod_{i \in I} \operatorname{Hom}_{\mathbf{A}}(\mathbf{M}, \mathbf{E}_i) \to \prod_{i \in I} \operatorname{Hom}_{\mathbf{A}}(\mathbf{M}', \mathbf{E}_i)\) to be surjective, it is necessary and sufficient that each of the homomorphisms \(\operatorname{Hom}_{\mathbf{A}}(\mathbf{M}, \mathbf{E}_i) \to \operatorname{Hom}_{\mathbf{A}}(\mathbf{M}', \mathbf{E}_i)\) be so (II, p.~10, prop. 5); this proves the proposition since \(\prod_{i \in I} \operatorname{Hom}_{\mathbf{A}}(\mathbf{M}, \mathbf{E}_i)\) is canonically identified with \(\operatorname{Hom}_{\mathbf{A}}(\mathbf{M}, \mathbf{E})\) .

PROPOSITION 10. --- Let E be an A-module. For E to be injective, it is necessary and sufficient that, for every ideal a of A and every A-homomorphism \(f: \mathfrak{a} \to \mathbb{E}\) there exists \(e \in \mathbb{E}\) such that \(f(a) = ae\) for all \(a \in \mathfrak{a}\) .

Suppose E is injective; let a be an ideal of A, \(f : a \to E\) an A-homomorphism, and denote by \(i : a \to A\) the canonical injection. Then the map

\[
\operatorname{Hom} _ {\mathbf {A}} (i, 1): \operatorname{Hom} _ {\mathbf {A}} (\mathrm{A}, \mathrm{E}) \rightarrow \operatorname{Hom} _ {\mathbf {A}} (\mathfrak {a}, \mathrm{E})
\]

is surjective (def. 2); if \(g \in \operatorname{Hom}_{\mathbf{A}}(\mathbf{A}, \mathbf{E})\) is such that \(f = g \circ i\) , we have

\[
f (a) = g (a) = a g (1)
\]

for all \(a \in \mathfrak{a}\) .

Conversely, suppose the condition in the statement holds; let M be an A-module, N a submodule of M, \(u: \mathbf{N} \to \mathbf{E}\) an A-homomorphism, and let us prove that there exists an A-homomorphism \(\overline{u}: \mathbf{M} \to \mathbf{E}\) extending \(u\) (cf. Lemma 3). Let \(\mathcal{P}\) be the set of pairs \((\mathbf{P}, v)\) where P is a submodule of M containing N and \(v\) a homomorphism from P into E extending \(u\) . The set \(\mathcal{P}\) ordered by the extension relation is inductive: if \((\mathbf{P}_{j}, v_{j})\) is a totally ordered family of elements of P, set \(Q = \cup P_{j}\) and let \(w : Q \to E\) be the unique map inducing \(v_{j}\) over \(P_{j}\) for all j; then \((\mathbf{Q}, w) \in \mathcal{P}\) and \((\mathbf{Q}, w)\) majorizes \((\mathbf{P}_{j}, v_{j})\) for all j. Then let \((\mathbf{P}, v)\) be a maximal element of P (E, III, p.~20, th. 2); it suffices to prove that P = M. Let \(x \in M\) and let a be the ideal of the \(a \in A\) such that \(ax \in P\) ; set \(f(a) = v(ax)\) for \(a \in a\) ; one thus obtains an A-homomorphism \(f : a \to E\) . Let then e be an element of E such that \(f(a) = ae\) for all \(a \in a\) . Set \(P' = P + Ax\) and let \(v' : P' \to E\) be the unique A-homomorphism such that \(v'(p + ax) = v(p) + ae\) for \(p \in P, a \in A\) ; then \((\mathbf{P}', v')\) belongs to P and majorizes \((\mathbf{P}, v)\) , hence \(P' = P\) , that is, \(x \in P\) , which completes the proof.

COROLLARY 1. --- If the ring A is left Noetherian, every module that is a direct sum of injective A-modules is injective.

Let \((\mathrm{E}_{i})_{i\in\mathbb{I}}\) be a family of injective A-modules, let E be their direct sum, a an ideal of A and \(u:a\to E\) an A-homomorphism. Since A is Noetherian, a is finitely generated, and consequently the canonical map

\[
\varphi : \bigoplus_ {i \in I} \operatorname{Hom} _ {\mathbf {A}} (\mathfrak {a}, E _ {i}) \rightarrow \operatorname{Hom} _ {\mathbf {A}} (\mathfrak {a}, E)
\]

is bijective; let \((u_{i})\) be the inverse image of u under \(\varphi\) . Since each \(E_{i}\) is injective, and the family \((u_{i})\) has finite support, there exists an element \((e_{i})_{i\in\mathbf{I}}\) of E such that \(u_{i}(a)=ae_{i}\) for all \(a\in\mathfrak{a}\) and for every \(i\in I\) , hence \(u(a)=a((e_{i}))\) for all \(a\in\mathfrak{a}\) , and E is injective.

Remark. --- If every A-module that is a direct sum of injective A-modules is injective, the ring A is left Noetherian (X, p.~170, Exercise 21).

Suppose A is an integral domain. One says that the A-module E is divisible if the homothety \(a_{E}\) is surjective for every nonzero element \(a\) of A.

COROLLARY 2. --- Suppose A is an integral domain.

\begin{enumerate}
\def\labelenumi{\alph{enumi})}
\item
  Every injective A-module is divisible.
\item
  Every torsion-free (II, p.~115) and divisible A-module is injective.
\item
  If A is principal, every divisible A-module is injective.
\end{enumerate}

If \(a \in A\) is nonzero, then \(a_{A}\) is injective; on the other hand, for every A-module E, the homothety \(a_{E}\) is canonically identified with

\[
\operatorname{Hom} \left(a _ {\mathrm{A}}, 1\right): \operatorname{Hom} _ {\mathrm{A}} (\mathrm{A}, \mathrm{E}) \rightarrow \operatorname{Hom} _ {\mathrm{A}} (\mathrm{A}, \mathrm{E}),
\]

Therefore, E is divisible if and only if Hom \((a_{\mathbf{E}}, 1_{\mathbf{E}})\) is surjective for all \(a \in \mathbf{A}\) nonzero. The assertion \(a)\) therefore follows from Definition 2 (X, p.~15).

Let E be a divisible A-module; suppose A is principal (resp. E is torsion-free) and let us prove that E is injective by applying prop. 10. Let a be an ideal of A and \(f: \mathfrak{a} \to \mathbf{E}\) an A-homomorphism. Let \(x \in \mathfrak{a}\) such that \(\mathfrak{a} = \mathbf{A}x\) (resp. such that \(x \neq 0\) if \(\mathfrak{a} \neq 0\) ), and let \(e \in \mathbf{E}\) such that \(xe = f(x)\) . Let us prove that for every \(a \in \mathfrak{a}\) , we have

\[
f (a) = a e;
\]

this is clear if \(a \in Ax\) , whence the assertion in the case where A is principal; if E is torsion-free and if \(xa \in \mathfrak{a}\) , we have \(xf(a) = f(ax) = axe\) , hence \(f(a) = ae\) since \(x\) is nonzero if \(a \neq 0\) .

Examples. --- 1) If A is an integral domain, the field of fractions K of A is an injective A-module. If A is principal, K/A is an injective A-module.

\begin{enumerate}
\def\labelenumi{\arabic{enumi})}
\setcounter{enumi}{1}
\item
  For example, the Z-modules Q and Q/Z are injective.
\item
  Let A be a principal ideal ring and let \(\iota\) be a nonzero element of A. Then A/aA is an injective A/aA-module (X, p.~170, exercise 20).
\end{enumerate}

